% 第 20 章 Graphical User Interface
\bichapter{图形用户界面}{Graphical User Interface}
\label{chap:gui}

PrimeTime 图形用户界面（GUI）允许您可视化设计数据，并使用直方图、原理图、抽象时钟图、波形图与数据表等分析工具查看分析结果。要学习如何使用 PrimeTime GUI，请参阅以下主题：
\begin{itemize}
  \item 打开与关闭 GUI
  \item 快速入门指南
  \item GUI 窗口
  \item 分析时序路径集合
  \item 在抽象时钟图中检查时钟路径
  \item 分析时钟域与时钟间关系
  \item 在原理图中检查时序路径与设计逻辑
  \item 查看与修改原理图
  \item 检查路径元素
  \item 查看对象属性
  \item 重新计算时序路径
  \item 布局视图与 GUI 中的 ECO
  \item 时序路径的交互式多场景分析
  \item GUI Tcl 命令
\end{itemize}

\section{打开与关闭 GUI}
\subsection*{Opening and Closing the GUI}

PrimeTime GUI 是在 Linux 上运行于 X 窗口环境中的窗口与菜单驱动界面。您可在启动 PrimeTime 时打开 GUI，也可在 PrimeTime 会话期间打开。

打开 PrimeTime GUI 时，工具会读取 GUI 设置与首选项文件，并打开新的 PrimeTime 主窗口：
\begin{itemize}
  \item 设置文件执行基本设置任务，例如初始化变量与声明设计库；首选项文件设置原理图与抽象时钟图视图属性以及全局应用首选项。
  \item 主窗口包含用于执行时序分析与其他分析任务的菜单、工具栏与视图窗口。
\end{itemize}

您可在会话期间随时打开或关闭 GUI。例如，当需要在 \cmd{pt\_shell} 中执行耗时任务或批处理时，可关闭 GUI；完成后再重新打开 GUI 进行可视化分析。

\begin{seeAlsoBox}
\begin{itemize}
  \item 打开 GUI
  \item 关闭 GUI
\end{itemize}
\end{seeAlsoBox}

\subsection{打开 GUI}
\subsubsection*{Opening the GUI}

打开 GUI 前，请确认 DISPLAY 环境变量已设置为您的 Linux 显示名称。也可在启动会话时同时设置 DISPLAY 变量。

要启动 PrimeTime 会话并打开 GUI，输入：
\begin{lstlisting}
% pt_shell -gui
\end{lstlisting}

若要在启动 PrimeTime 会话时设置 DISPLAY 环境变量，可使用 \opt{-display} 选项，其中 host\_name 为 Linux 显示终端名称。例如：
\begin{lstlisting}
% pt_shell -gui -display 192.180.50.155:0.0
\end{lstlisting}

要在 PrimeTime 会话期间打开 GUI，在 \cmd{pt\_shell} 提示符下输入：
\begin{lstlisting}
pt_shell> gui_start
\end{lstlisting}

\subsection{关闭 GUI}
\subsubsection*{Closing the GUI}

关闭 GUI 时，设计仍保留在内存中，命令行提示符在 shell 中保持活动。

要在不退出 pt_shell 的情况下关闭 GUI：
选择菜单 File > Close GUI；或在控制台或 \cmd{pt\_shell} 命令行输入 \cmd{gui\_stop}。

要关闭 GUI 并完全退出 PrimeTime 会话，选择 File > Exit。

\section{快速入门指南}
\subsection*{Quick Start Guide}

以下步骤介绍 GUI 的部分功能。请从已加载、已约束且已完成时序更新的设计开始。
在 GUI 中创建要分析的路径集合。例如：
\begin{lstlisting}
pt_shell> set mypaths [get_timing_paths -nworst 20 -max_paths 500 \
  -slack_lesser_than 0.2]
\end{lstlisting}

打开 GUI：
\begin{lstlisting}
pt_shell> gui_start
\end{lstlisting}

\begin{enumerate}
  \item 单击 Path Collections 选项卡。
  \item 右键单击路径集合名称，选择 Show Paths。将显示路径 slack 直方图与路径集合中的路径列表。
  \item 在列表中选择并高亮感兴趣的路径（例如最差 slack 的路径），然后单击 Inspector 按钮打开路径检查器。
  \item 在 Path Inspector 中查看 Delay Profiles、Path Summary、Path Elements 与 Timing report 等选项卡。
  \item 在 Schematic 视图中查看路径逻辑；使用 View Settings 面板调整显示。
  \item 使用 Clock > Clock Analyzer 分析时钟域关系。
\end{enumerate}

\figplaceholder{Figure 302}{PrimeTime 主窗口}{fig:pt-main-window}

\section{GUI 窗口}
\subsection*{GUI Windows}

PrimeTime GUI 在窗口中显示信息，为执行不同任务提供灵活工作环境。这些窗口在 X 窗口环境中独立运行，但共享内存中的同一设计、当前时序信息以及工具中的全局选择数据。

打开 GUI 时自动显示 PrimeTime 主窗口，提供多种分析视图。每个应用窗口包含标题栏、菜单、工具栏与面板、状态栏，以及显示特定设计信息的视图窗口。视图窗口与面板出现在工具栏与状态栏之间的工作区。

\figplaceholder{Figure 302}{PrimeTime 主窗口结构}{fig:gui-main}

\subsection{视图窗口}
\subsubsection*{View Windows}

视图窗口是 GUI 窗口内显示设计信息的子窗口。打开新视图窗口时，GUI 在工作区底部显示对应选项卡。

在视图窗口内单击可激活该视图（标题栏高亮）。重叠时活动视图位于最前。可通过选项卡或 Window 菜单激活视图；也可使用 Window > Next Window 或 Previous 循环切换。

含多个子视图的窗口为每个视图提供选项卡；打开时显示默认视图，单击其他选项卡可切换。

\subsection{视图设置面板}
\subsubsection*{View Settings Panel}

可通过 View Settings 面板查看并修改活动视图的显示设置。显示或隐藏面板：View > Toolbars > View Settings，或按 F8。可用设置取决于窗口类型。

默认情况下，更改选项后选项与 Apply 按钮变蓝；单击 Apply 使更改生效。可通过 Options 菜单或齿轮图标操作面板；可保存/恢复首选项或写入脚本。Auto apply 选项可跳过 Apply 按钮立即应用所有更改。

\figplaceholder{Figure 303}{视图设置面板}{fig:view-settings}

\subsection{工具栏与面板}
\subsubsection*{Toolbars and Panels}

工具栏为视图窗口中最常用命令与交互操作提供快捷访问；面板是增强工具栏，包含设置选项或处理设计数据的工具。

显示或隐藏工具栏/面板：View > Toolbars > toolbar\_name。

主窗口工具栏包括：Clock、View Zoom、History of Zoom/Pan、Schematics、View Tools、Highlight、ECO 等。

\figplaceholder{Figure 304}{主窗口工具栏}{fig:main-toolbars}

\subsection{层次浏览器}
\subsubsection*{Hierarchy Browser}

层次浏览器用于遍历设计层次并选择设计部分以进行后续分析。若未打开层次浏览器视图，选择 View > Hierarchy Browser 打开。

左窗格用 [+] 与 [–] 浏览单元层次；选择左窗格单元后，右窗格显示该单元内容。可选择列出对象类型：层次单元、所有单元、pin/port、net 等。

可用上下箭头键滚动对象列表；列宽不足时悬停指针可显示 InfoTip。单击列标题可排序，可调整列宽。

在层次浏览器中选择对象时，其他视图（如原理图）中也会同步选中。

\figplaceholder{Figure 305}{层次浏览器}{fig:hier-browser}

\subsection{控制台}
\subsubsection*{Console}

控制台用于执行 pt_shell 命令、查看文本格式 PrimeTime 响应以及命令历史。底部 pt_shell 按钮右侧的长框中可输入命令。

控制台操作与 pt_shell 终端窗口输入输出一致；可在终端或控制台任一处输入命令并查看响应。不需要控制台时可关闭，仅使用终端。

控制台提供终端没有的视图，底部选项卡可选择显示文本类型：
\begin{itemize}
  \item Log 选项卡：显示工具 pt_shell 响应。
  \item History 选项卡：查看最近执行命令历史。
\end{itemize}

右下角齿轮图标提供搜索、选择、复制与重用控制台内容的选项。

\figplaceholder{Figure 306}{控制台}{fig:console}

\subsection{对象选择}
\subsubsection*{Object Selection}

可在视图窗口中单击对象，或通过 Select 菜单命令选择对象以进行后续操作。交互式多选时按住 Ctrl 键；否则每次新选择会取消先前选择。

所选对象在所有视图中同步选中，不仅限于当前视图。例如某单元出现在两个原理图窗口与层次浏览器列表中，在任一处选中该单元，三处均会高亮。

\subsection{选择时序路径}
\subsubsection*{Selecting Timing Paths}

可选择特定时序路径，或选择设计中 slack 最差的一条或多条路径。例如显示路径 profile 前需先选择路径。简便方法是在 Path Collections 表中单击感兴趣路径。

创建路径集合示例：
\begin{lstlisting}
set mypath1 [get_timing_paths -from ... -to ...]
\end{lstlisting}

\subsection{设置 GUI 首选项}
\subsubsection*{Setting GUI Preferences}

打开 GUI 时从首选项文件加载 GUI 首选项。默认系统首选项针对大多数设计已优化。

设置 GUI 首选项：View > Preferences。更改后 GUI 自动将新设置保存到主目录下 \cmd{.synopsys\_pt\_prefs.tcl}；下次打开 GUI 时从此文件加载。

\section{分析时序路径集合}
\subsection*{Analyzing Timing Path Collections}

可分析时序路径集合以确定设计中时序失败位置。路径分析器（path analyzer）基于可用属性规则对路径分类；可选择预定义分类规则或定义自定义规则，也可添加子类别。

路径分析器是高级时序分析工具，支持自定义趋势分析。例如：
\begin{itemize}
  \item 按 path group 分类路径，查看违例是否集中于特定 path group。
  \item 为 slack 违例超过一个时钟周期的路径创建自定义规则，判断是否需要 multicycle path 约束。
  \item 按 start clock 与 end clock 分类，识别跨域路径（两值不同）。
  \item 为物理分区等块打 block mark，再按 block mark 属性分类，查看失败是否来自特定块。
\end{itemize}

打开路径分析器：Timing > Path Analyzer。

\figplaceholder{Figure 307}{初始路径分析器}{fig:path-analyzer-init}

进一步分析需加载时序路径集合（Collections 按钮）并对路径分类（Create 按钮）。

\figplaceholder{Figure 308}{加载并分类后的路径分析器}{fig:path-analyzer-cat}

\subsection{加载路径集合}
\subsubsection*{Loading Path Collections}

要将时序路径集合加载到路径分析器：
\begin{enumerate}
  \item 在路径分析器中单击 Collections 按钮。
  \item 在对话框中选择要加载的集合，单击 OK。
\end{enumerate}

\subsection{对时序路径分类}
\subsubsection*{Categorizing the Timing Paths}

要对时序路径分类：
1. 单击 Create 按钮。
2. 选择预定义规则（如 End Point、Start Clock、End Clock、Path Group 等）或自定义规则。
3. 单击 OK。

要删除类别或子类别，在 Categories 树中右键选择 Remove。

\subsection{将路径类别保存到文件}
\subsubsection*{Saving Path Categories to a File}

要将路径类别保存为脚本文件：在路径分析器中选择类别，右键选择 Save Categories，指定文件名。

\subsection{从文件加载路径类别}
\subsubsection*{Loading Path Categories From a File}

要从脚本文件加载路径类别：右键选择 Load Categories，选择先前保存的文件。

\subsection{标记块并应用块分类规则}
\subsubsection*{Marking Blocks and Applying Block Category Rules}

可为设计中的块（如物理分区）打 block mark，然后使用基于 block mark 属性的分类规则分析路径。

预定义块分类规则示例：
\begin{table}[htbp]
\centering
\caption{块的分类规则}
\small
\begin{tabular}{|l|l|l|}
\hline
\tblhead{规则名} & \tblhead{属性} & \tblhead{定义} \\
\hline
Inside block & block\_mark & 路径完全位于带 mark 的块内 \\
Crossing block & block\_mark & 路径穿越块边界 \\
Starting in block & block\_mark & 路径起点在块内 \\
Ending in block & block\_mark & 路径终点在块内 \\
\hline
\end{tabular}
\end{table}

\subsection{创建自定义分类规则}
\subsubsection*{Creating Custom Category Rules}

创建自定义分类规则：
1. 单击 Create，选择 Custom Rule。
2. 输入规则名称。
3. 在 Category Attribute Chooser 中选择属性。
4. 可选：在 Filter Attribute Chooser 中定义过滤表达式。
5. 单击 OK。

使用斜杠（/）分隔属性值可创建层次化子类别。

\section{在抽象时钟图中检查时钟路径}
\subsection*{Examining Clock Paths in an Abstract Clock Graph}

抽象时钟图（abstract clock graph）以图形方式显示时钟树结构与传播。可展开/折叠分支、显示/隐藏时钟元素、查看时钟延迟随时间布局等。

\figplaceholder{Figure 316}{抽象时钟图窗口}{fig:abstract-clk-graph}

\subsection{打开抽象时钟图视图}
\subsubsection*{Opening Abstract Clock Graph Views}

创建显示所有时钟的抽象时钟图：Clock > Abstract Clock Graph > All Clocks。

创建显示选定时钟的图：先选择时钟对象，再选择 Clock > Abstract Clock Graph > Selected Clocks。

\subsection{显示与隐藏时钟路径元素}
\subsubsection*{Displaying and Hiding Clock Path Elements}

在 View Settings 面板 Elements 选项卡中，可选择显示或隐藏：时钟源、缓冲器、反相器、ICG 单元、寄存器时钟 pin、端口等。

展开弧或 metacell：选中后右键 Expand Selected，或双击。

\figplaceholder{Figure 317}{展开选定弧}{fig:expand-arc}

\subsection{撤销与重做更改}
\subsubsection*{Reversing and Reapplying Changes}

在抽象时钟图视图中可撤销/重做操作（Edit > Undo / Redo，或工具栏按钮）。

\subsection{折叠选定侧分支}
\subsubsection*{Collapsing Selected Side Branches}

可折叠侧分支为单行以腾出空间：
1. 选择要折叠的分支根。
2. 右键 Collapse Branches。

展开已折叠侧分支：选择分支根，右键 Expand Branches。

\figplaceholder{Figure 318}{分支折叠示例}{fig:branch-collapse}

\subsection{查看时钟延迟随时间变化}
\subsubsection*{Viewing Clock Latency Over Time}

可在抽象时钟图视图中按调度关系显示时钟延迟，分析 skew、瓶颈、插入延迟与 ICG 内在路径延迟。

操作步骤：
1. 在 View Settings 面板单击 Settings 选项卡。
2. 选择 Show latency。
3. 单击 Apply。

x 轴表示相对时间的时钟树逻辑；逻辑门按输入时钟 pin 延迟左对齐。
\figplaceholder{Figure 319}{抽象时钟图中的时钟延迟布局}{fig:clk-latency-layout}

\subsection{层次化显示时钟树}
\subsubsection*{Displaying the Clock Trees Hierarchically}

默认抽象时钟图以扁平方式显示每条时钟路径。可在 View Settings > Elements 中选择 Hierarchy，将图折叠到顶层设计，再展开各层次单元。可选择逻辑层次（默认）或物理层次（Physical hierarchy）。

\section{分析时钟域与时钟间关系}
\subsection*{Analyzing Clock Domains and Clock-to-Clock Relationships}

使用时钟分析器（clock analyzer）识别设计中时钟间关系并驱动时钟分析，帮助确定哪些时钟域相互通信并识别时钟域交叉（CDC）。

时钟域由主时钟或生成时钟及其派生时钟组成；生成时钟有独立子域。打开时钟分析器：Clock > Clock Analyzer。

窗口左侧为层次化时钟树视图，右侧为时钟矩阵视图。
\figplaceholder{Figure 320}{时钟分析器}{fig:clock-analyzer}

时钟树视图显示主时钟名称；选择主时钟即选择整个时钟域（含其下生成时钟域）。时钟矩阵中每行每列代表一个时钟，用编号标注。

矩阵单元表示时钟域或特定 launch-capture 时钟对关系；字母与颜色标识路径约束及同步/异步关系。折叠时钟域时单元显示约束摘要（浅棕色单元）。

选中矩阵单元后，可在 Legend、Launch Clock、Capture Clock 框查看信息。右键可对 launch-capture 对执行 Analyze failing paths、Report exceptions、Report timing、Report clock skew 等操作。

可将时钟树或矩阵数据导出为 CSV 文件。

\subsubsection{查询 Launch-Capture 时钟对}
\paragraph*{Querying Launch-Capture Clock Pairs}

使用 Query 工具查询矩阵单元：
1. 单击时钟分析器窗口右上角 Query 按钮。
2. 单击矩阵单元，Query 面板显示 launch/capture 时钟与 interclock 关系。

\subsubsection{选择时钟或时钟域}
\paragraph*{Selecting Clocks or Clock Domains}

可在时钟树或矩阵视图中单击选择时钟；Ctrl+单击多选；Shift+单击选范围。
对给定主时钟深度选择所有生成时钟：选中 launch 时钟，右键 Deep Select all Generated Clocks。

选中时钟后可右键创建 Clock Graph of Selected Clocks、Clock Graph for Clock Domain、Schematic of Selected Clocks、Select Source Pins or Ports 等。

\subsubsection{排序时钟树视图}
\paragraph*{Sorting the Clock Tree View}

单击列标题降序排序，再次单击升序；层次结构保持。

\subsubsection{按名称查找时钟}
\paragraph*{Finding Clocks by Name}

选中时钟树中时钟，右键 Find Clock(s) 显示 Find 工具栏。可按名称、功能定义（如路径上 pin）等查找。

\figplaceholder{Figure 321}{Find 工具栏}{fig:find-toolbar}

\subsubsection{过滤时钟域}
\paragraph*{Filtering Clock Domains}

可使用过滤条件缩小显示的时钟集合，Focus 于特定时钟域分析。

\subsubsection{保存时钟属性或矩阵约束数据}
\paragraph*{Saving Clock Attribute or Clock Matrix Constraint Data}

可将时钟树属性或 launch-capture 矩阵约束导出为 CSV。

\subsubsection{时钟矩阵符号与颜色}
\paragraph*{Clock Matrix Symbols and Colors}

矩阵单元使用字母标识关系类型，例如：
\begin{itemize}
  \item S — 同步（synchronous）
  \item A — 异步（asynchronous）
  \item E — 显式路径（explicit path）
  \item I — 隐式路径（implicit path）
  \item 等（详见原书颜色图例）
\end{itemize}

\section{在原理图中检查时序路径与设计逻辑}
\subsection*{Examining Timing Paths and Design Logic in a Schematic}

实例原理图（instance schematic）可图形显示时序路径、选定设计逻辑或 fanin/fanout 逻辑，用于可视化分析时序与逻辑并指导其他分析任务。

实例原理图在扁平单页原理图中显示单元实例、pin、net、port、总线、ripper 与层次穿越，可跨多个层次级别。实例可为块（层次单元）或叶单元。

打开实例原理图：
1. 选择要在原理图中显示的设计对象或时序路径（层次浏览器、路径表、Select > Paths From/To/Through、Select > Fanin/Fanout 等）。
2. 选择 Schematic > Schematic View。

\figplaceholder{Figure 325}{实例原理图}{fig:instance-schematic}

含多层次时可用彩色边界按层次组织对象，可下移/上移层次查看块内容。可将缓冲器/反相器链或未重要块折叠为抽象 metacell。可在活动原理图视图中增删逻辑（不改变网表，不影响其他视图）。

\subsection{检查层次单元}
\subsubsection*{Examining Hierarchical Cells}

默认以扁平单页显示可跨层次的时序路径与逻辑；层次单元初始显示为折叠 metacell。
展开所有层次：Schematic > Expand > All Hierarchy。
展开选定层次：选中后 Schematic > Expand > Selected Objects 或双击。
折叠所有层次：Schematic > Collapse > All Hierarchy。
折叠选定层次：Schematic > Collapse > Selected Hierarchy By Parent。

\figplaceholder{Figure 326}{层次展开与折叠视图}{fig:hier-expand-collapse}

\subsection{在设计层次中下移或上移}
\subsubsection*{Moving Down or Up the Design Hierarchy}

下移：选中层次单元，Schematic > Move Down 或双击。
上移：Schematic > Move Up。

\begin{noteBox}
下移/上移时 GUI 会更新原理图选项卡上显示的设计名称。
\end{noteBox}

\subsection{显示或隐藏缓冲器与反相器}
\subsubsection*{Displaying or Hiding Buffers and Inverters}

默认扁平显示含层次穿越（菱形）的时序路径。可折叠缓冲器/反相器链或树为 metacell：
\begin{itemize}
  \item Schematic > Collapse > All Buffers/Inverters/Crossings By Chain
  \item Schematic > Collapse > All Buffers/Inverters/Crossings By Tree
  \item 选中对象后 Schematic > Collapse > Selected Buffers/Inverters/Crossings
\end{itemize}

\subsection{展开或折叠总线}
\subsubsection*{Expanding or Collapsing Buses}

可展开总线 net 与 pin 查看各位，或折叠为总线符号。

\section{查看与修改原理图}
\subsection*{Viewing and Modifying Schematics}

原理图视图提供缩放、平移、高亮、标注与查询工具。View Settings 面板可配置显示颜色、线宽、pin 标签、路径高亮样式等。

Select 菜单支持按名称选择/高亮对象；Query 工具可查看 pin/net/单元属性。

可撤销/重做原理图中的展开、折叠、增删对象等操作（Edit > Undo/Redo）。

\section{检查路径元素}
\subsection*{Inspecting Timing Path Elements}

路径检查器（Path Inspector）提供时序路径的详细视图，包括：
\begin{itemize}
  \item Delay Profiles — 路径延迟分布
  \item Path Summary — 路径摘要（slack、组、时钟等）
  \item Path Elements — 路径上各 pin 的到达时间、转换时间、增量延迟
  \item Timing report — 文本格式时序报告
\end{itemize}

配置路径检查器：在 Path Inspector 窗口 Preferences 中设置显示列与格式。

可从路径分析器或路径表选中路径后单击 Inspector 打开。

\figplaceholder{Figure 330}{路径检查器}{fig:path-inspector}

\section{查看对象属性}
\subsection*{Viewing Object Attributes}

Attribute Browser 与 Query 面板可查看设计对象属性。Attribute Group Manager 可创建、修改、复制与重命名属性组。

\subsection{查看当前选择}
\subsubsection*{Viewing the Current Selection}

使用 \cmd{get\_selection} 或在 GUI 中查看选择摘要。Select > Selection Summary 可报告当前选中对象数量与类型。

\subsection{管理属性组}
\subsubsection*{Managing Attribute Groups}

打开 Attribute Group Manager：Tools > Attribute Group Manager。
可创建自定义属性组、从首选项恢复、修改 Basic/Placement/SchemAnnotAttr/Timing 组。

复制组：选中组后单击 Copy；重命名：Rename 按钮编辑名称。

\section{重新计算时序路径}
\subsection*{Recalculating Timing Paths}

可通过重新计算路径表访问并比较正常路径与其重新计算对应项。使用带 \opt{-pba\_mode path} 的 \cmd{get\_timing\_paths} 生成重新计算路径集合。例如：
\begin{lstlisting}
set my_recalc_path [get_timing_paths -from ... -to ... -pba_mode path]
\end{lstlisting}

\subsection{比较正常与重新计算路径}
\subsubsection*{Comparing Normal and Recalculated Paths}

在 GUI 中：从路径分析器选择路径，选择 Timing > Recalculated Path Comparison Table。表格显示 endpoint、path group、正常 slack 与重新计算 slack；违例单元红色，通过为绿色。

\figplaceholder{Figure 345}{重新计算路径表}{fig:recalc-path-table}

\subsection{比较正常与重新计算路径 pin}
\subsubsection*{Comparing Normal and Recalculated Path Pins}

选择重新计算路径表中的行，选择 Timing > Recalculated Path Pin Comparison Table。并排比较 pin 转换时间、路径延迟、增量延迟等；正负变化分别以绿/红高亮。

仅当正常与重新计算路径具有相同 through pin 数量时启用 pin 比较表。

\figplaceholder{Figure 346}{路径 pin 比较表}{fig:recalc-pin-table}

\section{布局视图与 GUI 中的 ECO}
\subsection*{Layout View and ECOs in the GUI}

要在 GUI 中查看芯片布局以准备工程变更指令（ECO），在主窗口选择 Window > Layout window。

使用布局视图需用 \cmd{set\_eco\_options} 指定物理数据（IC Compiler II 或 LEF/DEF 格式）。用 \cmd{check\_eco} 检查物理数据有效性。

可在 GUI 中生成并查看 ECO，包括插入缓冲器、删除缓冲器、调整单元尺寸；需要 PrimeTime-ADV-PLUS 许可证。更改后可轻松撤销。

\figplaceholder{Figure 347}{GUI 中的布局视图}{fig:layout-view}

\figplaceholder{Figure 348}{GUI 中的 ECO 菜单}{fig:eco-menu}

启用撤销：确保 ECO 菜单中 Enabled 选项有勾选；不使用时可禁用以节省内存。

可高亮时序路径并仅显示路径上实际 net 形状（含相关 via），不含无关 fanout。

\figplaceholder{Figure 349}{Net 形状显示}{fig:net-shape}

\figplaceholder{Figure 350}{Net 形状显示选项}{fig:net-shape-opts}

使用 ECO > Insert Buffer 或 ECO > Size Cell（\cmd{insert\_buffer} / \cmd{size\_cell}）执行手动 ECO 时，工具自动在内存中执行增量仅路径时序更新，可立即查看更改的时序影响。

布局期间 on-route 缓冲器引导会遵守 placement blockage。要禁用自动增量更新，单击 GUI 底部 Fast Timing Path update Enabled 旁的 Disable。

\figplaceholder{Figure 351}{GUI 中的手动 ECO 操作}{fig:manual-eco-gui}

在布局窗口显示彩色单元密度图与 pin 密度图：View > Map > Cell Density 或 Pin Density（需 PrimeTime-ADV-PLUS）。
显示功耗密度图：View > Map > Cell Power Density（需 PrimePower 分析数据）。

\figplaceholder{Figure 352}{布局视图中的单元密度图}{fig:cell-density-map}

\figplaceholder{Figure 353}{布局视图中的 pin 密度图}{fig:pin-density-map}

\figplaceholder{Figure 354}{布局视图中的单元功耗密度图}{fig:power-density-map}

\section{时序路径的交互式多场景分析}
\subsection*{Interactive Multi-Scenario Analysis of Timing Paths}

PrimeTime GUI 提供交互式多场景分析（IMSA）模式，可从不同运行或场景恢复时序路径集合并一起分析，作为跨分析运行调试与跟踪违例的快速工具。

示例流程：
\begin{lstlisting}
\section{PrimeTime Session 1}
...
set my_paths1 [get_timing_paths ...]
save_session -only_timing_paths $my_paths1 /dir/s1

\section{PrimeTime Session 2}
...
set my_paths2 [get_timing_paths ...]
save_session -only_timing_paths $my_paths2 /dir/s2

\section{PrimeTime multi-scenario analysis session}
gui_start
restore_session /dir/s1
restore_session /dir/s2
\section{在 GUI 中显示、分类、排序多会话路径 ...}
\end{lstlisting}

\begin{noteBox}
恢复的仅路径会话须具有相同网表与时钟（约束可不同），且 save/restore 会话的 PrimeTime 版本须匹配。IMSA 会话中命令集受限，仅可运行分析已恢复路径集合的命令。
\end{noteBox}

\subsection{交互式多场景分析流程}
\subsubsection*{Interactive Multi-Scenario Analysis Flow}

IMSA 流程可：加载并调试多组时序路径集合；分类与排序（见路径分析器）；按指定值平移类别 slack 以便组间比较；增删信息列。

使用步骤：
1. 用 get_timing_paths 创建路径集合。
2. 用 save_session -only_timing_paths 保存。
3. 在新会话中 gui_start，restore_session 加载各会话路径。
4. 在路径分析器 Setup Collections 中选择集合并 Apply。
5. 创建分类规则（如 End Point）分析路径。

\subsubsection{平移类别的 Slack}
\paragraph*{Shifting the Slack for a Category}

在 IMSA 流程中可将类别 slack 平移用户定义值：
1. 在路径分析器中选择类别。
2. 右键 Shift Category，在 Shift Histogram 对话框设置平移值与父类别。
3. 在 Table Histogram 中选择 Slack+Shift 或 Slack 列查看直方图。

\figplaceholder{Figure 355}{Shift Histogram 对话框}{fig:shift-histogram}

\figplaceholder{Figure 356}{Table Histogram 对话框}{fig:table-histogram}

\figplaceholder{Figure 357}{路径 Slack 与 Slack+Shift}{fig:slack-shift}

\subsubsection{IMSA 保存与恢复的属性}
\paragraph*{IMSA Attributes Saved and Restored}

save_session -only_timing_paths 保存路径对象及 IMSA 相关属性，包括 session 名称、场景标识、约束差异标记等，便于跨会话比较。

\section{GUI Tcl 命令}
\subsection*{GUI Tcl Commands}

PrimeTime 支持一组控制与查询 GUI 窗口的 Tcl 命令，可用于编写脚本控制 GUI 外观与内容。例如：
\begin{itemize}
  \item 创建菜单命令
  \item 创建工具栏
  \item 添加命令按钮
  \item 向 Favorites 面板添加按钮
  \item 设置对象颜色
  \item 在布局窗口中添加颜色高亮与标注
\end{itemize}

GUI 控制能力类似于 IC Compiler II。列出 GUI 相关命令：
\begin{lstlisting}
pt_shell> help gui_*
 gui_change_highlight  # 更改全局高亮状态
 gui_create_attrgroup    # 创建 GUI 属性组
 gui_create_pref_category # 创建首选项类别
 gui_create_pref_key     # 创建首选项键值对
 gui_create_vm           # 创建新 Visual Mode
 ...
\end{lstlisting}

\begin{lstlisting}
pt_shell> man gui_change_highlight
\end{lstlisting}

\cmd{gui\_change\_highlight} 用于操纵全局高亮对象集。

\subsection{GUI 菜单结构概览}
\subsubsection*{GUI Menu Structure Overview}

PrimeTime GUI 主菜单包括：
\begin{itemize}
  \item \textbf{File} — 打开/关闭 GUI、保存/恢复会话、退出
  \item \textbf{Edit} — 撤销/重做、复制、首选项
  \item \textbf{View} — 工具栏、视图窗口、布局图、密度图、缩放
  \item \textbf{Select} — 按名称/类型/路径选择对象
  \item \textbf{Timing} — Path Analyzer、Recalculated Path 比较、时序报告
  \item \textbf{Clock} — Clock Analyzer、Abstract Clock Graph
  \item \textbf{Schematic} — 原理图视图、展开/折叠、层次导航
  \item \textbf{ECO} — 插入/删除缓冲器、调整单元尺寸（需布局视图）
  \item \textbf{Window} — 窗口管理、Layout window
  \item \textbf{Help} — 帮助与文档
\end{itemize}

\subsection{路径分析器高级功能}
\subsubsection*{Advanced Path Analyzer Features}

路径分析器支持多列排序、过滤表达式、导出 CSV、与 Path Inspector 联动、按 session 分组（IMSA）等。

预定义分类规则包括：Startpoint、Endpoint、Start Clock、End Clock、Path Group、Logic Levels、Fanout、Net Length 等。

\subsubsection{过滤表达式语法}
\paragraph*{Filter Expression Syntax}

自定义过滤支持比较运算符（==、!=、>、<）与逻辑组合（\texttt{\&\&}、\texttt{||}）。属性名区分大小写，字符串值用引号括起。

\subsection{原理图视图工具}
\subsubsection*{Schematic View Tools}

\begin{itemize}
  \item Zoom In/Out/Fit — 缩放与适应窗口
  \item Pan — 平移视图
  \item Highlight — 高亮选中对象
  \item Query — 查询 pin/cell/net 属性
  \item Add Fanin/Fanout — 添加扇入/扇出逻辑
  \item Add Path — 添加时序路径到原理图
\end{itemize}

\subsection{抽象时钟图高级操作}
\subsubsection*{Advanced Abstract Clock Graph Operations}

支持按时钟域过滤、显示/隐藏 ICG、显示生成时钟关系、导出时钟树报告、与 Clock Analyzer 联动跳转。

\figplaceholder{Figure 322}{时钟矩阵符号图例}{fig:clk-matrix-legend}

\subsection{波形与数据表视图}
\subsubsection*{Waveform and Data Table Views}

GUI 支持将路径延迟数据以波形图与表格形式显示，便于分析路径上各段增量延迟贡献。

\subsection{GUI 会话保存与恢复}
\subsubsection*{GUI Session Save and Restore}

save_session 可保存完整会话或仅路径集合；restore_session 恢复后可在 GUI 中继续分析。IMSA 模式依赖 save/restore 跨场景比较路径。

\begin{lstlisting}
save_session -only_timing_paths $paths /path/to/session
restore_session /path/to/session
\end{lstlisting}

\subsection{GUI 与 pt_shell 协同}
\subsubsection*{GUI and pt_shell Interaction}

\begin{noteBox}
GUI 与 pt_shell 共享同一设计数据库；在任一处执行的命令会更新双方视图。耗时 update_timing 建议在 pt_shell 终端执行，完成后在 GUI 刷新视图。
\end{noteBox}

\subsection{GUI 相关环境变量}
\subsubsection*{GUI Environment Variables}

\begin{itemize}
  \item DISPLAY — X11 显示（必需）
  \item SYNOPSYS_PT_GUI_PREF — 首选项文件路径覆盖
  \item TCL_LIBRARY — Tcl 脚本库路径
\end{itemize}

\subsection{gui\_* 命令完整列表（节选）}
\subsubsection*{gui\_* Command Reference (Selected)}

\begin{longtable}{|L{0.38\textwidth}|L{0.55\textwidth}|}
\hline
\tblhead{命令} & \tblhead{功能} \\
\hline
\texttt{gui\_start} / \texttt{gui\_stop} & 打开/关闭 GUI \\
\texttt{gui\_change\_highlight} & 操纵全局高亮对象集 \\
\texttt{gui\_create\_attrgroup} & 创建属性组 \\
\texttt{gui\_create\_pref\_category} & 创建首选项类别 \\
\texttt{gui\_create\_pref\_key} & 创建首选项键值 \\
\texttt{gui\_create\_vm} & 创建 Visual Mode \\
\texttt{gui\_execute\_menu\_item} & 执行菜单项 \\
\texttt{gui\_get\_current\_window} & 获取当前窗口 \\
\texttt{gui\_set\_attrgroup\_value} & 设置属性组值 \\
\texttt{gui\_set\_pref\_key\_value} & 设置首选项值 \\
\texttt{gui\_show\_window} & 显示指定窗口 \\
\texttt{gui\_update\_display} & 刷新显示 \\
\hline
\end{longtable}

